% GENERATED FILE. Edit canonical lessons, then run npm run generate:books.
% canonical-source-hash: fnv1a64:0d5d445722243e52
% canonical-lessons: ML-C465-purappetal, ML-C465-varavu, ML-C465-kalatamasam, ML-C465-kattunilkkuka, ML-C465-otikkayaruka

\chapter{Departures and Arrivals}
\label{ch:ml-departures-and-arrivals}
\hlchaptermodality{465}

\begin{chapteropening}
\textbf{By the end of this chapter:} \emph{I can say departure, arrival, a delay, to wait standing and to run and board.}

Complete the last lesson of chapter 465: I can say departure, arrival, a delay, to wait standing and to run and board.
\end{chapteropening}

\section[\texorpdfstring{\ml{പുറപ്പെടൽ} (puṟappeṭal)}{പുറപ്പെടൽ (puṟappeṭal)}]{\ml{പുറപ്പെടൽ} (puṟappeṭal) — departure}

\label{lesson:ML-C465-purappetal}



\begin{quote}
Before the new one: say the Malayalam for an inspection, then the Malayalam for a foreigner.
\end{quote}

\subsection*{You'll want to know: \ml{പുറപ്പെടൽ}}
\textbf{\ml{പുറപ്പെടൽ}} — \emph{puṟappeṭal} — \textquotedblleft{}departure\textquotedblright{}.

\begin{quote}
\ml{വിമാനത്തിന്റെ} \ml{പുറപ്പെടൽ} \ml{രാത്രി} \ml{പത്ത്} \ml{മണിക്കാണ്}.
\end{quote}

\begin{quote}
\emph{vimānattinṟe puṟappeṭal rātri pattŭ maṇikkāṇŭ.}
\end{quote}

\begin{quote}
The flight's departure is at ten at night.
\end{quote}

The noun of \ml{പുറപ്പെടുക}, to set out: the ending -\ml{അൽ} turns a verb into a noun.

\begin{cousinweb}
\textbf{\ml{പുറപ്പെടൽ}} (\emph{puṟappeṭal}) is native Malayalam, the verbal noun of \ml{പുറപ്പെടുക}, to set out, which is built on \ml{പുറം}, the outside (DEDR 4333), with the ending -\ml{അൽ}.
\end{cousinweb}

\begin{grammarlens}[title={What you've built}]
One more word for getting around.
\end{grammarlens}

\subsection*{Your turn}
\begin{itemize}
\raggedright
  \item \emph{Say it:} \emph{puṟappeṭal}
  \item \emph{Say it:} \emph{puṟappeṭal}, once more
  \item \emph{Say it:} \emph{vimānattinṟe puṟappeṭal rātri pattŭ maṇikkāṇŭ.} — The flight's departure is at ten at night.
  \item \emph{Recall:} say the Malayalam for an inspection, then the Malayalam for a foreigner, then say \emph{puṟappeṭal} again
\end{itemize}

\begin{tcolorbox}[breakable,colback=teal!4,colframe=teal!35!black,title={Before you move on}]
What does \ml{പുറപ്പെടൽ} mean? (\textquotedblleft{}Departure\textquotedblright{}.) Say it once more.
\end{tcolorbox}
\section[\texorpdfstring{\ml{വരവ്} (varavŭ)}{വരവ് (varavŭ)}]{\ml{വരവ്} (varavŭ) — arrival; coming}

\label{lesson:ML-C465-varavu}



\begin{quote}
Before the new one: say the Malayalam for a foreigner, then the Malayalam for departure.
\end{quote}

\subsection*{You'll want to know: \ml{വരവ്}}
\textbf{\ml{വരവ്}} — \emph{varavŭ} — \textquotedblleft{}arrival; coming\textquotedblright{}.

\begin{quote}
\ml{അമ്മയുടെ} \ml{വരവിനായി} \ml{കുട്ടികൾ} \ml{കാത്തിരുന്നു}.
\end{quote}

\begin{quote}
\emph{ammayuṭe varavināyi kuṭṭikaḷ kāttirunnu.}
\end{quote}

\begin{quote}
The children waited for their mother's arrival.
\end{quote}

The noun of \ml{വരുക}, to come. Its opposite is \ml{പോക്ക്}, a going.

\begin{cousinweb}
\textbf{\ml{വരവ്}} (\emph{varavŭ}) is native Dravidian, the noun of \ml{വരുക}, to come (DEDR 5270).
\end{cousinweb}

\begin{grammarlens}[title={What you've built}]
One more word for getting around.
\end{grammarlens}

\subsection*{Your turn}
\begin{itemize}
\raggedright
  \item \emph{Say it:} \emph{varavŭ}
  \item \emph{Say it:} \emph{varavŭ}, once more
  \item \emph{Say it:} \emph{ammayuṭe varavināyi kuṭṭikaḷ kāttirunnu.} — The children waited for their mother's arrival.
  \item \emph{Recall:} say the Malayalam for a foreigner, then the Malayalam for departure, then say \emph{varavŭ} again
\end{itemize}

\begin{tcolorbox}[breakable,colback=teal!4,colframe=teal!35!black,title={Before you move on}]
What does \ml{വരവ്} mean? (\textquotedblleft{}Arrival; coming\textquotedblright{}.) Say it once more.
\end{tcolorbox}
\section[\texorpdfstring{\ml{കാലതാമസം} (kālatāmasaṁ)}{കാലതാമസം (kālatāmasaṁ)}]{\ml{കാലതാമസം} (kālatāmasaṁ) — a delay}

\label{lesson:ML-C465-kalatamasam}



\begin{quote}
Before the new one: say the Malayalam for departure, then the Malayalam for arrival.
\end{quote}

\subsection*{You'll want to know: \ml{കാലതാമസം}}
\textbf{\ml{കാലതാമസം}} — \emph{kālatāmasaṁ} — \textquotedblleft{}a delay\textquotedblright{}.

\begin{quote}
\ml{മഴ} \ml{കാരണം} \ml{തീവണ്ടിക്ക്} \ml{രണ്ട്} \ml{മണിക്കൂർ} \ml{കാലതാമസം} \ml{ഉണ്ടായി}.
\end{quote}

\begin{quote}
\emph{maḻa kāraṇaṁ tīvaṇṭikkŭ raṇṭŭ maṇikkūr kālatāmasaṁ uṇṭāyi.}
\end{quote}

\begin{quote}
Because of the rain the train had a two-hour delay.
\end{quote}

A formal word, common in announcements; in speech people say \ml{വൈകി}, it is late.

\begin{cousinweb}
\textbf{\ml{കാലതാമസം}} (\emph{kālatāmasaṁ}) is a Malayalam compound built from \ml{കാലം}, time, and \ml{താമസം}, delay, both borrowed from Sanskrit: \emph{kāla}, time, and \emph{tāmasa}, \textquotedblleft{}dark, sluggish\textquotedblright{} (Monier-Williams).
\end{cousinweb}

\begin{grammarlens}[title={What you've built}]
One more word for getting around.
\end{grammarlens}

\subsection*{Your turn}
\begin{itemize}
\raggedright
  \item \emph{Say it:} \emph{kālatāmasaṁ}
  \item \emph{Say it:} \emph{kālatāmasaṁ}, once more
  \item \emph{Say it:} \emph{maḻa kāraṇaṁ tīvaṇṭikkŭ raṇṭŭ maṇikkūr kālatāmasaṁ uṇṭāyi.} — Because of the rain the train had a two-hour delay.
  \item \emph{Recall:} say the Malayalam for departure, then the Malayalam for arrival, then say \emph{kālatāmasaṁ} again
\end{itemize}

\begin{tcolorbox}[breakable,colback=teal!4,colframe=teal!35!black,title={Before you move on}]
What does \ml{കാലതാമസം} mean? (\textquotedblleft{}A delay\textquotedblright{}.) Say it once more.
\end{tcolorbox}
\section[\texorpdfstring{\ml{കാത്തുനിൽക്കുക} (kāttunilkkuka)}{കാത്തുനിൽക്കുക (kāttunilkkuka)}]{\ml{കാത്തുനിൽക്കുക} (kāttunilkkuka) — to wait standing}

\label{lesson:ML-C465-kattunilkkuka}



\begin{quote}
Before the new one: say the Malayalam for arrival, then the Malayalam for a delay.
\end{quote}

\subsection*{You'll want to know: \ml{കാത്തുനിൽക്കുക}}
\textbf{\ml{കാത്തുനിൽക്കുക}} — \emph{kāttunilkkuka} — \textquotedblleft{}to wait standing\textquotedblright{}.

\begin{quote}
\ml{മഴയത്ത്} \ml{ഞങ്ങൾ} \ml{ബസിനായി} \ml{അര} \ml{മണിക്കൂർ} \ml{കാത്തുനിന്നു}.
\end{quote}

\begin{quote}
\emph{maḻayattŭ ñaṅṅaḷ basināyi ara maṇikkūr kāttuninnu.}
\end{quote}

\begin{quote}
We stood waiting for the bus in the rain for half an hour.
\end{quote}

Past: \ml{കാത്തുനിന്നു}, \emph{kāttuninnu}.

\ml{കാത്തിരിക്കുക} is to wait in general, often sitting; \ml{കാത്തുനിൽക്കുക} is to wait on your feet, at a stop or a gate.

\begin{cousinweb}
\textbf{\ml{കാത്തുനിൽക്കുക}} (\emph{kāttunilkkuka}) is a native Dravidian compound of \ml{കാക്കുക}, to keep watch or wait (DEDR 1416), and \ml{നിൽക്കുക}, to stand (DEDR 3675).
\end{cousinweb}

\begin{grammarlens}[title={What you've built}]
One more word for getting around.
\end{grammarlens}

\subsection*{Your turn}
\begin{itemize}
\raggedright
  \item \emph{Say it:} \emph{kāttunilkkuka}
  \item \emph{Say it:} \emph{kāttunilkkuka}, once more
  \item \emph{Say it:} \emph{maḻayattŭ ñaṅṅaḷ basināyi ara maṇikkūr kāttuninnu.} — We stood waiting for the bus in the rain for half an hour.
  \item \emph{Recall:} say the Malayalam for arrival, then the Malayalam for a delay, then say \emph{kāttunilkkuka} again
\end{itemize}

\begin{tcolorbox}[breakable,colback=teal!4,colframe=teal!35!black,title={Before you move on}]
What does \ml{കാത്തുനിൽക്കുക} mean? (\textquotedblleft{}To wait standing\textquotedblright{}.) Say it once more.
\end{tcolorbox}
\section[\texorpdfstring{\ml{ഓടിക്കയറുക} (ōṭikkayaṟuka)}{ഓടിക്കയറുക (ōṭikkayaṟuka)}]{\ml{ഓടിക്കയറുക} (ōṭikkayaṟuka) — to run and board}

\label{lesson:ML-C465-otikkayaruka}



\begin{quote}
Before the new one: say the Malayalam for a delay, then the Malayalam for to wait standing.
\end{quote}

\subsection*{You'll want to know: \ml{ഓടിക്കയറുക}}
\textbf{\ml{ഓടിക്കയറുക}} — \emph{ōṭikkayaṟuka} — \textquotedblleft{}to run and board\textquotedblright{}.

\begin{quote}
\ml{തീവണ്ടി} \ml{നീങ്ങുന്നതിന്} \ml{മുൻപ്} \ml{അവൻ} \ml{ഓടിക്കയറി}.
\end{quote}

\begin{quote}
\emph{tīvaṇṭi nīṅṅunnatinŭ munpŭ avan ōṭikkayaṟi.}
\end{quote}

\begin{quote}
He ran and got on before the train moved off.
\end{quote}

Past: \ml{ഓടിക്കയറി}, \emph{ōṭikkayaṟi}.

\ml{ഓടി}, having run, plus \ml{കയറുക}, to get in. The vehicle takes -\ml{ഇൽ}: \ml{ബസിൽ} \ml{ഓടിക്കയറി}.

\begin{cousinweb}
\textbf{\ml{ഓടിക്കയറുക}} (\emph{ōṭikkayaṟuka}) is a native Dravidian compound of \ml{ഓടുക}, to run (DEDR 1041), and \ml{കയറുക}, to climb in (DEDR 1255).
\end{cousinweb}

\begin{grammarlens}[title={What you've built}]
One more word for getting around.
\end{grammarlens}

\subsection*{Your turn}
\begin{itemize}
\raggedright
  \item \emph{Say it:} \emph{ōṭikkayaṟuka}
  \item \emph{Say it:} \emph{ōṭikkayaṟuka}, once more
  \item \emph{Say it:} \emph{tīvaṇṭi nīṅṅunnatinŭ munpŭ avan ōṭikkayaṟi.} — He ran and got on before the train moved off.
  \item \emph{Recall:} say the Malayalam for a delay, then the Malayalam for to wait standing, then say \emph{ōṭikkayaṟuka} again
\end{itemize}

\begin{tcolorbox}[breakable,colback=teal!4,colframe=teal!35!black,title={Before you move on}]
What does \ml{ഓടിക്കയറുക} mean? (\textquotedblleft{}To run and board\textquotedblright{}.) Say it once more.
\end{tcolorbox}
